\manualchapter{Control reference}{sec:controls}
Panel map: \textbf{1} pitch, period and pitch/FM inputs;
\textbf{2} FM and noise mode;
\textbf{3} level and level CV;
\textbf{4} audio outputs.

Three tone columns are followed by the noise column. Tones have frequency
and FM controls; noise has period and LFSR controls. Every column has level
CV and an audio output.

\begin{tabularx}{\textwidth}{@{}l l X@{}}
\toprule Control & Range; default & Operation \\\midrule
Tone frequency & $\pm2.5$ oct.; 0 & C4 reference; V/OCT adds octaves. \\
FM depth & $-1$--$+1$; 0 & Adds $aV/5$ octaves to the tone. \\
Noise period & 0--3; 0 & Adds $V/2$, floors/clamps, then reverses the register. \\
Level & 0--15; 7 & Multiplies by $V/10$, rounds and clamps. \\
LFSR & Off/on; off & Chooses the noise sequence; CV changes the switch state. \\
\bottomrule
\end{tabularx}

Noise setting 0 uses Tone 3's period. Settings 1--3 select fixed chip periods
of 1024, 512 and 256 clock units, respectively. Each 2\,V of period CV advances
one setting. The LFSR CV has 2\,V high and 0.01\,V low thresholds.

Tone pitch and FM normal forward from Tone 1 to 3; their first inputs default
to 0 and 5\,V. Level normals through all four columns from an initial 10\,V.
FM is exponential and its depth knob also serves as fine tuning when empty.
The chip's frequency and amplitude registers introduce discrete steps;
they are part of the oscillator response.
